%% ma: L12-B06-0036
%% lop: 12
%% chuong: 2
%% bai: 6
%% dang: 12.6.15
%% loai: TLN
%% muc: TH
%% dap_an: "600"
%% nguon: "LaTeX giáo viên gửi 21/10/2026 (M05, Vectơ và hệ trục tọa độ trong không gian, 107 câu), Phần III Câu 5"
%% kien_thuc: "Bài 6 (tích vô hướng, ứng dụng tính công)"
%% trang_thai: chinh-thuc
%% ghi_chu: "Giáo viên đã duyệt 21/10/2026. Hình gốc vẽ lực F nghiêng 30° nhưng ghi 60°; đã vẽ lại đúng góc 60°"
\begin{ex}
Một thùng hàng được kéo trượt trên mặt sàn nằm ngang bằng một lực $\overrightarrow{F}$ có độ lớn $150$ N. Sợi dây kéo hợp với mặt sàn một góc $60^\circ$ (tham khảo hình bên dưới). Thùng hàng dịch chuyển được một đoạn đường dài $8$ m. Tính công do lực kéo sinh ra (đơn vị: J).
\begin{center}
\begin{tikzpicture}
  \fill[gray!20] (-1,0) rectangle (1,1.2); \draw (-1,0) rectangle (1,1.2);
  \draw[thick] (-2,0)--(5,0);
  \draw[->,thick] (1,.6)--++(60:2.4) node[above left]{$\overrightarrow{F}$};
  \draw[phu] (1,.6)--(3.4,.6);
  \draw (1.8,.6) arc (0:60:.8); \node at (2.15,1.05) {$60^\circ$};
  \draw[->] (-1,-.35)--(3,-.35) node[midway,below]{$\overrightarrow{d}$, $|\overrightarrow{d}|=8$ m};
\end{tikzpicture}
\end{center}
\shortans{600}
\loigiai{$A=\overrightarrow{F}\cdot\vec{d}=150\cdot8\cdot\cos60^\circ=600$ (J).}
\end{ex}
